\scanpage{15}
% archive sheet number: 14
\noindent\begin{tabular}[t]{@{}c@{}}
U N I V E R S I T E\\
DES SCIENCES ET TECHNIQUES\\
DU LANGUEDOC\\
\textasciicircum\textasciicircum\textasciicircum\textasciicircum\textasciicircum\textasciicircum\textasciicircum\textasciicircum\textasciicircum\textasciicircum\textasciicircum\textasciicircum\\[3pt]
Service de la Scolarité\\
JP/GJ\\[3pt]
\uline{Diffusion} : Liste A\\
N\textdegree{}2417
\end{tabular}\hfill\begin{tabular}[t]{@{}l@{}}
MONTPELLIER, le 19 AVRIL 1977\\
\end{tabular}

\begin{center}
\uline{AVIS DE SOUTENANCE D'UNE THESE DE DOCTORAT DE}\\[3pt]
\uline{SPECIALITE : CHIMIE Mention CHIMIE-PHYSIQUE}
\end{center}

\uline{Nom et prénom du candidat} : BERLANDIER Gaston

\uline{Titre de la thèse} : Etude comparée des spectres de vibration de
\[
\mathrm{CH}_{3}\text{-}\mathrm{C}\hphantom{\mathrm{CC}}\mathrm{Cl}_{2}\text{-}\mathrm{CH}_{3},\quad
\mathrm{CD}_{3}\text{-}\mathrm{C}\hphantom{\mathrm{C}}\mathrm{Cl}_{2}\text{-}\mathrm{CD}_{3}
\quad\text{et}\quad\mathrm{CD}_{3}\text{-}\mathrm{CCl}_{2}\text{-}\mathrm{CH}_{3}
\]

\uline{Date, heure et lieu de soutenance prévus} : 23 AVRIL 1977 à 10 heures, Salle de Cours n\textdegree{} 10 - U.S.T.L.

\uline{Nom du Directeur de thèse} : Messieurs BRIBES et GAUFRES

\uline{Désignation du laboratoire où la thèse est préparée} : Laboratoire de Spectroscopie Moléculaire de l'Université des Sciences et Techniques du Languedoc.

\uline{Analyse de l'ouvrage} : Après avoir décrit la préparation des composés deutérés, ainsi que le contrôle de leur pureté isotopique, l'auteur présente et analyse les spectres de vibration des trois espèces, à savoir les spectres d'absorption et de diffusion des phases liquide et gazeuse. L'identification des fréquences fondamentales repose sur les formes des enveloppes des bandes en infrarouge comme en Raman, ainsi que sur le facteur de dépolarisation en diffusion. La recherche de la meilleure attribution compte tenu de la règle du produit et des informations expérimentales est faite systématiquement par ordinateur. 73 fréquences fondamentales sur 81 sont données avec une bonne certitude.

\uline{Durée approximative des recherches ayant abouti à la thèse} : Deux ans

\uline{Composition du jury proposé} :
\begin{center}
\begin{tabular}{@{}lll@{}}
MM. : & R. & GAUFRES\\
 & J. & PETRISSANS\\
 & J.L. & BRIBES\\
 & P. & SAUMAGNE\\
 & V. & TABACIK
\end{tabular}
\end{center}

\begin{flushright}
Le Secrétaire Général,\\
\handwritten{\unsure{E S|E Siau}}\\
\uline{E. SIAU}\\
\end{flushright}
